\chapter{Flux de développement}
\label{chap-dev-workflow}

\renewcommand{\headrulewidth}{0.3pt}
\fancyhead[LE]{\thepage} \fancyhead[RE]{\nouppercase{\textbf{\leftmark}}}
\fancyhead[RO]{\thepage} \fancyhead[LO]{\nouppercase{\textbf{\rightmark}}}
\fancyfoot[L,C,R]{}

\setlength{\parindent}{0.35cm}

\minitoc
% ================================================================

Ce chapitre rassemble les \textbf{règles et processus} qu'un développeur doit suivre en travaillant sur \cwv~: comment étendre l'application avec une nouvelle opération, la liste de contrôle à parcourir avant de committer, les conventions du projet, et comment le versionnage et la documentation du code sont gérés.

\section{Ajouter une nouvelle opération de dialogue}
\label{sec-dev-add-op}

Pour ajouter un nouveau flux au niveau dialogue, suivez trois étapes :

\begin{enumerate}
    \item définissez une dataclass \script{*Result} dans \script{ht/client/api\_types.py} ;
    \item implémentez \script{run\_<name>} dans \script{ht/client/operations\_client.py} --- il possède l'orchestration \script{fetch}~$\rightarrow$~\script{transform}~$\rightarrow$~\script{render} ;
    \item ajoutez une méthode de façade d'une ligne sur \script{HydrologicalTwinClient}, puis appelez-la depuis \script{ExploreData}.
\end{enumerate}

Les corps de méthode sur la classe client sont volontairement minces ($\leq$~5 instructions). Gardez-les ainsi.

\section{Lors d'une modification --- liste de contrôle}
\label{sec-dev-checklist}

\begin{itemize}
    \item vous touchez au code d'une boîte de dialogue ? $\rightarrow$ \script{frontend/Dialog*.py}. Passez par \script{ExploreData} ;
    \item vous ajoutez un nouveau flux au niveau dialogue ? $\rightarrow$ la recette en trois étapes ci-dessus (\cref{sec-dev-add-op}) ;
    \item vous ajoutez une nouvelle primitive de calcul ? $\rightarrow$ \script{ht/developer/handlers.py} (calcul) ou \script{services/public/twin\_io.py} (lectures d'état L3 / E/S disque) ; routez \textit{via} \script{dispatch.py} avec un nouveau \script{kind=} ;
    \item vous avez trouvé un import \script{qgis} / \script{PyQt5} dans \script{external/HydrologicalTwinAlphaSeries/} ? $\rightarrow$ bug. Déplacez-le vers \script{interface/} ou \script{frontend/} ;
    \item vous avez trouvé un module sous \script{services/} qui importe depuis \script{ht/client/} ou \script{ht/developer/} ? $\rightarrow$ bug. C'est un import \emph{ascendant} ; les imports doivent pointer uniquement vers le bas ;
    \item vous avez trouvé un module frontend qui contourne \script{ExploreData} pour atteindre les internes de HTAS ? $\rightarrow$ bug. Passez par la porte client ;
    \item vous avez trouvé une manipulation de données dans \script{ht/client/} (construction de \script{DataFrame}, remise en forme/agrégation de tableaux, mise en forme d'unités/chemins, E/S disque, maths du domaine) ? $\rightarrow$ bug. L1 ne fait qu'orchestrer.
\end{itemize}

\section{Conventions}
\label{sec-dev-conventions}

\begin{itemize}
    \item \textbf{Environnement :} géré par \script{pixi} (\script{pixi.toml}, \script{pixi.lock}). QGIS, Python et toutes les dépendances sont figés dans un unique environnement reproductible (\cref{sec-install-pixi}) ;
    \item \textbf{Aides à la synchronisation :} \script{./verifState.sh} (diagnostics) et \script{./updateProject.sh} (synchronisation ordonnée multi-dépôts / sous-modules) ;
    \item \textbf{Découverte pré-cycle-de-vie} (sans jumeau requis) : \script{HydrologicalTwinClient.detect\_from\_out\_caw(...)} et \script{detect\_project\_neighbors(...)} --- aides statiques utilisées par \script{DialogSettings} pour pré-remplir les chemins avant \script{configure} / \script{load} ;
    \item \textbf{Cale de compatibilité :} \script{cawaqsviz/backend}, si présent, ré-exporte les symboles HTAS. L'implémentation réelle ne vit que dans \script{external/HydrologicalTwinAlphaSeries/}.
\end{itemize}

\section{Mise à jour de version}

La version du plug-in est spécifiée dans le fichier \script{metadata.txt} figurant dans l'extrait de code \ref{code_metadata}.
\begin{lstlisting}[language=json, caption={Informations obligatoires à renseigner dans le fichier \script{metadata.txt}}, label={code_metadata}]
[general]
name=CaWaQS-Viz
qgisMinimumVersion=3.0
description=Description
version=0.1.16
author=Lise-Marie GIROD (Mines Paris - PSL), Nicolas Flipo (Mines Paris - PSL)
email=lise-marie.girod@minesparis.psl.eu

about=Python QGIS-based vizualization software for CaWaQS3.x

repository=https://gitlab.com/gviz/cawaqsviz
\end{lstlisting}

\section{Documentation du code}

La documentation du code est publiée à l'adresse suivante via GitLab Pages : \href{https://cawaqsviz-gviz-adf347ff6d9d45c174db091497a2dfc0f14305cdb4ed9070.gitlab.io/}{https://cawaqsviz-gviz-adf347ff6d9d45c174db091497a2dfc0f14305cdb4ed9070.gitlab.io/}. La documentation est maintenue à travers les docstrings des modules de l'application, en suivant le formatage montré dans l'exemple ci-dessous.


\begin{lstlisting}[language=Python, caption={Formatage des docstrings pour l'intégration automatique dans la documentation du code}, label={code_format_docstrings}]
class Exemple:
    """
    Resume des fonctionnalites de la classe.
    """

    def __init__(self, arg1):
        """
        Methode constructeur.

        :param arg1: descriptif du parametre
        :type arg1: str
        """
        self.arg1 = arg1

    def do_something(self, arg1: list) -> str:
        """
        Descriptif de la fonction.

        :param arg1: descriptif du parametre
        :type arg1: list
        :return: descriptif de ce que la fonction retourne
        :rtype: str
        """
        return ", ".join(map(str, arg1))
\end{lstlisting}
\vspace{1em}
La compilation est réalisée par \script{Sphinx}\footnote{dont la documentation est accessible via \href{https://www.sphinx-doc.org/en/master/}{https://www.sphinx-doc.org/en/master/}}. Le répertoire \script{docs/sources} contient les fichiers de configuration (\script{Conf.py}), ainsi que le contenu de la documentation. Lors de l'intégration de nouveaux modules dans le plug-in, leur ajout à l'index (\script{index.rst}) est nécessaire pour que la docstring du code apparaisse sur la page de documentation. Cet index intègre 3 fichiers : \script{frontend.rst}, \script{interface.rst}, \script{backend.rst}, qui contiennent les classes à documenter. L'intégration d'un nouveau module dans l'application requiert son ajout à ces fichiers au format de l'extrait de code \ref{code_format_rst}. Les mises à jour des fichiers sources de la documentation doivent être committées et synchronisées avec les dépôts.

\begin{lstlisting}[language=Python,caption={Formatage des docstrings pour l'intégration automatique dans la documentation du code}, label={code_format_rst}]
Module
------

.. automodule:: cawaqsviz.Module
   :members:
   :show-inheritance:
   :undoc-members:
\end{lstlisting}

\vspace{1em}

Enfin, à chaque nouveau commit sur les branches \script{Main} ou \script{Beta}, la documentation est compilée selon le contenu du répertoire \script{docs/source} du dépôt.

\begin{lstlisting}
image: python:3.9

stages:
  - build
  - deploy

before_script:
  - pip install -U pip
  - pip install -r requirements.txt

build-docs:
  stage: build
  script:
    - sphinx-build -b html docs/source public
  artifacts:
    paths:
      - public

pages:
  stage: deploy
  script:
    - echo "Pages deployed"
  artifacts:
    paths:
      - public
  only:
    - main
    - beta

\end{lstlisting}
